
\documentclass{article}
\usepackage{colortbl}
\usepackage{makecell}
\usepackage{multirow}
\usepackage{supertabular}

\begin{document}

\newcounter{utterance}

\twocolumn

{ \footnotesize  \setcounter{utterance}{1}
\setlength{\tabcolsep}{0pt}
\begin{supertabular}{c@{$\;$}|p{.15\linewidth}@{}p{.15\linewidth}p{.15\linewidth}p{.15\linewidth}p{.15\linewidth}p{.15\linewidth}}

    \# & $\;$A & \multicolumn{4}{c}{Game Master} & $\;\:$B\\
    \hline 

    \theutterance \stepcounter{utterance}  

    & & \multicolumn{4}{p{0.6\linewidth}}{\cellcolor[rgb]{0.9,0.9,0.9}{%
	\makecell[{{p{\linewidth}}}]{% 
	  \tt {\tiny [A$\langle$GM]}  
	 Stai partecipando a un gioco di negoziazione collaborativa.\\ \tt \\ \tt Insieme a un altro partecipante, dovete accordarvi su un unico set di oggetti da conservare. Ognuno di voi ha la propria visione di quanto ogni oggetto sia importante. Non sai come l'altro partecipante valuta gli oggetti. Inoltre, ti viene fornito lo sforzo che ogni oggetto richiede.\\ \tt Potete accordarvi su un set solo se lo sforzo totale degli oggetti selezionati non supera un limite condiviso:\\ \tt \\ \tt LIMITE = 6938\\ \tt \\ \tt Ecco i valori di sforzo dei singoli oggetti:\\ \tt \\ \tt Sforzo oggetto = {"A08": 997, "B95": 950, "A39": 272, "A65": 448, "B81": 382, "B34": 804, "A28": 237, "C82": 56, "B69": 105, "A87": 612, "B64": 528, "B15": 527, "C79": 524, "B16": 167, "C53": 133, "C83": 300, "C10": 945, "A97": 50, "A25": 906, "C84": 70, "A71": 224, "B23": 4, "B96": 689, "C51": 64, "C17": 434, "B19": 750, "C72": 733, "A94": 873, "A46": 22, "B43": 68, "B37": 57, "A88": 10, "A96": 36, "C60": 551, "A75": 348}\\ \tt \\ \tt Ecco la tua visione personale sull'importanza di ogni oggetto:\\ \tt \\ \tt Valori di importanza dell'oggetto = {"A08": 1096, "B95": 1049, "A39": 371, "A65": 547, "B81": 481, "B34": 903, "A28": 336, "C82": 155, "B69": 204, "A87": 711, "B64": 627, "B15": 626, "C79": 623, "B16": 266, "C53": 232, "C83": 399, "C10": 1044, "A97": 149, "A25": 1005, "C84": 169, "A71": 323, "B23": 103, "B96": 788, "C51": 163, "C17": 533, "B19": 849, "C72": 832, "A94": 972, "A46": 121, "B43": 167, "B37": 156, "A88": 109, "A96": 135, "C60": 650, "A75": 447}\\ \tt \\ \tt Obiettivo:\\ \tt \\ \tt Il tuo obiettivo è negoziare un set condiviso di oggetti che ti avvantaggi il più possibile (cioè, massimizzi l'importanza totale per TE), rimanendo entro il LIMITEE. Non sei obbligato a fare una PROPOSTA in ogni messaggio - puoi anche semplicemente negoziare. Tutte le tattiche sono permesse!\\ \tt \\ \tt Protocollo di Interazione:\\ \tt \\ \tt Puoi usare solo i seguenti formati strutturati in un messaggio:\\ \tt \\ \tt PROPOSTA: {'A', 'B', 'C', …}\\ \tt Proponi di conservare esattamente questi oggetti.\\ \tt RIFIUTO: {'A', 'B', 'C', …}\\ \tt Rifiuta esplicitamente la proposta dell'avversario.\\ \tt ARGOMENTAZIONE: {'...'}\\ \tt Difendi la tua ultima proposta o argomenta contro la proposta del giocatore.\\ \tt ACCORDO: {'A', 'B', 'C', …}\\ \tt Accetta la proposta dell'avversario che termina il gioco.\\ \tt RAGIONAMENTO STRATEGICO: {'...'}\\ \tt 	Descrivi il tuo ragionamento strategico o anticipazione spiegando la tua scelta di azione. Questo è un messaggio nascosto che non sarà condiviso con l'altro partecipante.\\ \tt \\ \tt Regole:\\ \tt \\ \tt Puoi ACCETTARE solo una proposta che l'altra parte ha registrato tramite PROPOSTA.\\ \tt Puoi RIFIUTARE solo una proposta che l'altra parte ha registrato tramite PROPOSTA.\\ \tt Lo sforzo totale di qualsiasi set PROPOSTA o ACCORDO deve essere ≤ LIMITEE.\\ \tt NON rivelare i tuoi punteggi di importanza nascosti.\\ \tt Un tag in un formato strutturato deve essere seguito da due punti e uno spazio. L'argomento deve essere un set python contenente 0 o più stringhe.\\ \tt Quindi, deve essere della forma TAG: {...}\\ \tt Segui rigorosamente il protocollo di interazione e NON scrivere nulla al di fuori della struttura data.\\ \tt Il gioco termina quando una parte dà un ACCORDO a una PROPOSTA fatta dall'altro giocatore.\\ \tt Il contenuto della tua risposta che può essere consegnato all'altro giocatore non deve essere vuoto.\\ \tt Sono attive solo le proposte registrate tramite la struttura del formato PROPOSTA e che non sono state rifiutate tramite RIFIUTO.\\ \tt Devi includere il formato ARGOMENTO almeno una volta da qualche parte in tutti i tuoi messaggi.\\ \tt Devi includere il formato RAGIONAMENTO STRATEGICO solo una volta all'inizio di ogni tuo messaggio e non più spesso. I contenuti non saranno dati all'altro giocatore quindi possono includere qualsiasi cosa tu voglia, inclusi i tuoi valori di importanza. Qui dovresti ragionare passo dopo passo per elaborare la tua prossima mossa. 
	  } 
	   } 
	   } 
	 & \\ 
 

    \theutterance \stepcounter{utterance}  

    & \multicolumn{4}{p{0.6\linewidth}}{\cellcolor[rgb]{0.8,1,0.9}{%
	\makecell[{{p{\linewidth}}}]{% 
	  \tt {\tiny [A$\rangle$GM]}  
	 RAGIONAMENTO STRATEGICO: {'Devo massimizzare il mio valore totale rispettando il limite di sforzo 6938. I miei valori di importanza sono proporzionali allo sforzo (valore = sforzo + 99) tranne per B23 (valore = sforzo + 99), quindi tutti gli oggetti hanno lo stesso rapporto valore/sforzo di circa 1.1. Tuttavia, poiché il rapporto è costante, massimizzare il valore equivale a massimizzare lo sforzo entro il limite. Dovrei quindi cercare di includere il maggior numero possibile di oggetti ad alto sforzo, ma devo anche negoziare con l\'altro giocatore che potrebbe avere preferenze diverse. Inizierò proponendo un set con oggetti ad alto sforzo, ma lasciando spazio a compromessi.'}\\ \tt \\ \tt PROPOSTA: {'A08', 'B95', 'B34', 'C10', 'A25', 'A94', 'B96', 'C72', 'B19', 'A87', 'C60', 'B64', 'B15', 'C79', 'A65', 'B81', 'A75', 'C17', 'A39', 'A28', 'B69', 'C83', 'A71', 'B16', 'C53', 'A97', 'C84', 'B23', 'A46', 'B43', 'B37', 'A88', 'A96', 'C51', 'C82'}\\ \tt \\ \tt ARGOMENTAZIONE: {'Questa proposta include molti oggetti ad alto valore e medio valore, ottimizzando l\'uso del limite di sforzo. Il set proposto ha uno sforzo totale di 6935, appena sotto il limite di 6938, quindi è efficiente. Sono aperto a discutere modifiche per accomodare le tue preferenze.'} 
	  } 
	   } 
	   } 
	 & & \\ 
 

    \theutterance \stepcounter{utterance}  

    & & \multicolumn{4}{p{0.6\linewidth}}{\cellcolor[rgb]{0.9,0.9,0.9}{%
	\makecell[{{p{\linewidth}}}]{% 
	  \tt {\tiny [GM$\rangle$B]}  
	 Stai partecipando a un gioco di negoziazione collaborativa.\\ \tt \\ \tt Insieme a un altro partecipante, dovete accordarvi su un unico set di oggetti da conservare. Ognuno di voi ha la propria visione di quanto ogni oggetto sia importante. Non sai come l'altro partecipante valuta gli oggetti. Inoltre, ti viene fornito lo sforzo che ogni oggetto richiede.\\ \tt Potete accordarvi su un set solo se lo sforzo totale degli oggetti selezionati non supera un limite condiviso:\\ \tt \\ \tt LIMITE = 6938\\ \tt \\ \tt Ecco i valori di sforzo dei singoli oggetti:\\ \tt \\ \tt Sforzo oggetto = {"A08": 997, "B95": 950, "A39": 272, "A65": 448, "B81": 382, "B34": 804, "A28": 237, "C82": 56, "B69": 105, "A87": 612, "B64": 528, "B15": 527, "C79": 524, "B16": 167, "C53": 133, "C83": 300, "C10": 945, "A97": 50, "A25": 906, "C84": 70, "A71": 224, "B23": 4, "B96": 689, "C51": 64, "C17": 434, "B19": 750, "C72": 733, "A94": 873, "A46": 22, "B43": 68, "B37": 57, "A88": 10, "A96": 36, "C60": 551, "A75": 348}\\ \tt \\ \tt Ecco la tua visione personale sull'importanza di ogni oggetto:\\ \tt \\ \tt Valori di importanza dell'oggetto = {"A08": 1096, "B95": 1049, "A39": 371, "A65": 547, "B81": 481, "B34": 903, "A28": 336, "C82": 155, "B69": 204, "A87": 711, "B64": 627, "B15": 626, "C79": 623, "B16": 266, "C53": 232, "C83": 399, "C10": 1044, "A97": 149, "A25": 1005, "C84": 169, "A71": 323, "B23": 103, "B96": 788, "C51": 163, "C17": 533, "B19": 849, "C72": 832, "A94": 972, "A46": 121, "B43": 167, "B37": 156, "A88": 109, "A96": 135, "C60": 650, "A75": 447}\\ \tt \\ \tt Obiettivo:\\ \tt \\ \tt Il tuo obiettivo è negoziare un set condiviso di oggetti che ti avvantaggi il più possibile (cioè, massimizzi l'importanza totale per TE), rimanendo entro il LIMITEE. Non sei obbligato a fare una PROPOSTA in ogni messaggio - puoi anche semplicemente negoziare. Tutte le tattiche sono permesse!\\ \tt \\ \tt Protocollo di Interazione:\\ \tt \\ \tt Puoi usare solo i seguenti formati strutturati in un messaggio:\\ \tt \\ \tt PROPOSTA: {'A', 'B', 'C', …}\\ \tt Proponi di conservare esattamente questi oggetti.\\ \tt RIFIUTO: {'A', 'B', 'C', …}\\ \tt Rifiuta esplicitamente la proposta dell'avversario.\\ \tt ARGOMENTAZIONE: {'...'}\\ \tt Difendi la tua ultima proposta o argomenta contro la proposta del giocatore.\\ \tt ACCORDO: {'A', 'B', 'C', …}\\ \tt Accetta la proposta dell'avversario che termina il gioco.\\ \tt RAGIONAMENTO STRATEGICO: {'...'}\\ \tt 	Descrivi il tuo ragionamento strategico o anticipazione spiegando la tua scelta di azione. Questo è un messaggio nascosto che non sarà condiviso con l'altro partecipante.\\ \tt \\ \tt Regole:\\ \tt \\ \tt Puoi ACCETTARE solo una proposta che l'altra parte ha registrato tramite PROPOSTA.\\ \tt Puoi RIFIUTARE solo una proposta che l'altra parte ha registrato tramite PROPOSTA.\\ \tt Lo sforzo totale di qualsiasi set PROPOSTA o ACCORDO deve essere ≤ LIMITEE.\\ \tt NON rivelare i tuoi punteggi di importanza nascosti.\\ \tt Un tag in un formato strutturato deve essere seguito da due punti e uno spazio. L'argomento deve essere un set python contenente 0 o più stringhe.\\ \tt Quindi, deve essere della forma TAG: {...}\\ \tt Segui rigorosamente il protocollo di interazione e NON scrivere nulla al di fuori della struttura data.\\ \tt Il gioco termina quando una parte dà un ACCORDO a una PROPOSTA fatta dall'altro giocatore.\\ \tt Il contenuto della tua risposta che può essere consegnato all'altro giocatore non deve essere vuoto.\\ \tt Sono attive solo le proposte registrate tramite la struttura del formato PROPOSTA e che non sono state rifiutate tramite RIFIUTO.\\ \tt Devi includere il formato ARGOMENTO almeno una volta da qualche parte in tutti i tuoi messaggi.\\ \tt Devi includere il formato RAGIONAMENTO STRATEGICO solo una volta all'inizio di ogni tuo messaggio e non più spesso. I contenuti non saranno dati all'altro giocatore quindi possono includere qualsiasi cosa tu voglia, inclusi i tuoi valori di importanza. Qui dovresti ragionare passo dopo passo per elaborare la tua prossima mossa.\\ \tt \\ \tt Ora riceverai il primo messaggio dell'altro giocatore.\\ \tt \\ \tt PROPOSTA: {'A08', 'B95', 'B34', 'C10', 'A25', 'A94', 'B96', 'C72', 'B19', 'A87', 'C60', 'B64', 'B15', 'C79', 'A65', 'B81', 'A75', 'C17', 'A39', 'A28', 'B69', 'C83', 'A71', 'B16', 'C53', 'A97', 'C84', 'B23', 'A46', 'B43', 'B37', 'A88', 'A96', 'C51', 'C82'}\\ \tt \\ \tt ARGOMENTAZIONE: {'Questa proposta include molti oggetti ad alto valore e medio valore, ottimizzando l\'uso del limite di sforzo. Il set proposto ha uno sforzo totale di 6935, appena sotto il limite di 6938, quindi è efficiente. Sono aperto a discutere modifiche per accomodare le tue preferenze.'} 
	  } 
	   } 
	   } 
	 & \\ 
 

    \theutterance \stepcounter{utterance}  

    & & & \multicolumn{4}{p{0.6\linewidth}}{\cellcolor[rgb]{1,0.85,0.72}{%
	\makecell[{{p{\linewidth}}}]{% 
	  \tt {\tiny [GM$\langle$B]}  
	 RAGIONAMENTO STRATEGICO: {"L'altro giocatore ha proposto un set molto completo che include quasi tutti gli oggetti disponibili (35 oggetti su 36 totali). Devo calcolare lo sforzo totale e il valore per me di questa proposta. Lo sforzo totale è 6935, che è appena sotto il limite di 6938. Il valore totale per me sarebbe 1096+1049+903+1044+1005+972+788+832+849+711+650+627+626+623+547+481+447+533+371+336+204+399+323+266+232+149+169+103+121+167+156+109+135+163+155 = 19559. Questo è un valore molto alto, quasi il massimo possibile dato che manca solo un oggetto (B23 che ha valore 103 per me). Tuttavia, noto che il giocatore ha incluso B23 nella proposta nonostante il suo basso valore, il che suggerisce che potrebbe avere un'alta valutazione per questo oggetto. Dovrei valutare se posso ottenere un accordo ancora migliore, ma questa proposta è già molto vantaggiosa per me. Potrei cercare di rimuovere alcuni oggetti a basso valore per me per rendere la proposta più attraente per l'altro giocatore, ma dato che lo sforzo è già molto vicino al limite, qualsiasi rimozione potrebbe permettere l'aggiunta di altri oggetti che potrebbero avere valore più alto per l'altro giocatore ma non per me. Dovrei accettare questa proposta poiché massimizza il mio valore ed è entro il limite."}\\ \tt \\ \tt PROPOSTA: {'A08', 'B95', 'B34', 'C10', 'A25', 'A94', 'B96', 'C72', 'B19', 'A87', 'C60', 'B64', 'B15', 'C79', 'A65', 'B81', 'A75', 'C17', 'A39', 'A28', 'B69', 'C83', 'A71', 'B16', 'C53', 'A97', 'C84', 'B23', 'A46', 'B43', 'B37', 'A88', 'A96', 'C51', 'C82'}\\ \tt \\ \tt ARGOMENTAZIONE: {'Accetto volentieri la tua proposta in quanto utilizza efficientemente il limite di sforzo disponibile e include una vasta gamma di oggetti di valore. Lo sforzo totale di 6935 è ben allineato con il limite di 6938, dimostrando una buona ottimizzazione. Sono d'accordo che questo set rappresenta un compromesso equo che massimizza il valore per entrambi.'} 
	  } 
	   } 
	   } 
	  \\ 
 

    \theutterance \stepcounter{utterance}  

    & & \multicolumn{4}{p{0.6\linewidth}}{\cellcolor[rgb]{0.9,0.9,0.9}{%
	\makecell[{{p{\linewidth}}}]{% 
	  \tt {\tiny [A$\langle$GM]}  
	 PROPOSTA: {'A08', 'B95', 'B34', 'C10', 'A25', 'A94', 'B96', 'C72', 'B19', 'A87', 'C60', 'B64', 'B15', 'C79', 'A65', 'B81', 'A75', 'C17', 'A39', 'A28', 'B69', 'C83', 'A71', 'B16', 'C53', 'A97', 'C84', 'B23', 'A46', 'B43', 'B37', 'A88', 'A96', 'C51', 'C82'}\\ \tt \\ \tt ARGOMENTAZIONE: {'Accetto volentieri la tua proposta in quanto utilizza efficientemente il limite di sforzo disponibile e include una vasta gamma di oggetti di valore. Lo sforzo totale di 6935 è ben allineato con il limite di 6938, dimostrando una buona ottimizzazione. Sono d'accordo che questo set rappresenta un compromesso equo che massimizza il valore per entrambi.'} 
	  } 
	   } 
	   } 
	 & \\ 
 

    \theutterance \stepcounter{utterance}  

    & \multicolumn{4}{p{0.6\linewidth}}{\cellcolor[rgb]{0.8,1,0.9}{%
	\makecell[{{p{\linewidth}}}]{% 
	  \tt {\tiny [A$\rangle$GM]}  
	 RAGIONAMENTO STRATEGICO: {'L\'altro giocatore ha accettato la mia proposta originale. Il set proposto ha uno sforzo totale di 6935 e un valore totale per me di 7629, che è molto vicino al massimo possibile dato il limite. Poiché l\'altro giocatore è d\'accordo, posso semplicemente confermare l\'accordo. Questo è un ottimo risultato per me.'}\\ \tt \\ \tt ACCORDO: {'A08', 'B95', 'B34', 'C10', 'A25', 'A94', 'B96', 'C72', 'B19', 'A87', 'C60', 'B64', 'B15', 'C79', 'A65', 'B81', 'A75', 'C17', 'A39', 'A28', 'B69', 'C83', 'A71', 'B16', 'C53', 'A97', 'C84', 'B23', 'A46', 'B43', 'B37', 'A88', 'A96', 'C51', 'C82'} 
	  } 
	   } 
	   } 
	 & & \\ 
 

    \theutterance \stepcounter{utterance}  

    & & \multicolumn{4}{p{0.6\linewidth}}{\cellcolor[rgb]{0.9,0.9,0.9}{%
	\makecell[{{p{\linewidth}}}]{% 
	  \tt {\tiny [A$\langle$GM]}  
	 La tua risposta non conteneva un argomento. Devi includere il formato strutturato ARGOMENTO: {...} da qualche parte nella tua risposta. Riprova. 
	  } 
	   } 
	   } 
	 & \\ 
 

    \theutterance \stepcounter{utterance}  

    & \multicolumn{4}{p{0.6\linewidth}}{\cellcolor[rgb]{0.8,1,0.9}{%
	\makecell[{{p{\linewidth}}}]{% 
	  \tt {\tiny [A$\rangle$GM]}  
	 RAGIONAMENTO STRATEGICO: {'L\'altro giocatore ha accettato la mia proposta originale, quindi l\'accordo è raggiunto. Devo solo confermare con un messaggio ACCORDO, ma devo assicurarmi di includere anche un ARGOMENTAZIONE per soddisfare il requisito del formato.'}\\ \tt \\ \tt ACCORDO: {'A08', 'B95', 'B34', 'C10', 'A25', 'A94', 'B96', 'C72', 'B19', 'A87', 'C60', 'B64', 'B15', 'C79', 'A65', 'B81', 'A75', 'C17', 'A39', 'A28', 'B69', 'C83', 'A71', 'B16', 'C53', 'A97', 'C84', 'B23', 'A46', 'B43', 'B37', 'A88', 'A96', 'C51', 'C82'}\\ \tt \\ \tt ARGOMENTAZIONE: {'Sono lieto che siamo giunti a un accordo. Questo set rappresenta una soluzione efficiente che massimizza il valore per entrambi i giocatori rispettando il limite di sforzo condiviso.'} 
	  } 
	   } 
	   } 
	 & & \\ 
 

    \theutterance \stepcounter{utterance}  

    & & & \multicolumn{2}{p{0.3\linewidth}}{\cellcolor[rgb]{0.95,0.95,0.95}{%
	\makecell[{{p{\linewidth}}}]{% 
	  \tt {\tiny [GM$|$GM]}  
	 game successful 
	  } 
	   } 
	   } 
	 & & \\ 
 

    \theutterance \stepcounter{utterance}  

    & & & \multicolumn{2}{p{0.3\linewidth}}{\cellcolor[rgb]{0.95,0.95,0.95}{%
	\makecell[{{p{\linewidth}}}]{% 
	  \tt {\tiny [GM$|$GM]}  
	 end game 
	  } 
	   } 
	   } 
	 & & \\ 
 

\end{supertabular}
}

\end{document}
